\section{The estimator and its properties}\label{sec:method}

This section states the construction of Sections~\ref{sec:intro}--\ref{sec:worked}
formally and establishes its properties. The order is: notation and the object being
decomposed (Section~\ref{sec:setup}); the population identity and the covariance form of
its remainder (Section~\ref{sec:transport}), which Section~\ref{sec:bregman} shows is not
specific to the quadratic score; what happens when the grouping variable is made coarser
(Section~\ref{sec:coarsening}) or is missing a relevant covariate altogether
(Section~\ref{sec:sensitivity}); the finite-sample identity and its $O(NK)$ computation
(Sections~\ref{sec:finite}--\ref{sec:efficient}); why the counterfactual leaves each case
out rather than fitting the group's own frequencies (Section~\ref{sec:attenuation}); and
inference, both finite-sample and asymptotic
(Sections~\ref{sec:inference}--\ref{sec:asymptotic}). Proofs are in
Appendix~\ref{app:proofs}. A reader who wants only what is needed to use the estimator can
read Sections~\ref{sec:setup}, \ref{sec:transport}, \ref{sec:finite}
and~\ref{sec:inference} and skip the rest.

\subsection{Setup: what is being decomposed}\label{sec:setup}

Let $q_0(\cdot\mid x)$ and $q_1(\cdot\mid x)$ be the predictive distributions of an
old and a new frozen model over $K$ labels. Let $S(q,y)$ be a bounded proper score,
oriented so that larger is better, and let $C=\phi(X)$ be a declared discrete
grouping variable. In the relation-prediction example of Section~\ref{sec:intro}, $C$ is the ordered
subject--object class pair; in a long-tailed classification benchmark it is a coarse
superclass. We call a value of $C$ a \emph{cell} and require it to be declared before the
audit is run. Define the
instance-specific \emph{gain vector}
\begin{equation}
H_x(y)=S(q_1(\cdot\mid x),y)-S(q_0(\cdot\mid x),y),\qquad y=1,\ldots,K.
\label{eq:gain-vector}
\end{equation}
Unlike the realized gain $H_X(Y)$, this vector can be evaluated at a counterfactual
label without rerunning either model.

Our default is the quadratic score
\begin{equation}
S_{\rm B}(q,y)=2q_y-\lVert q\rVert_2^2,
\label{eq:brier}
\end{equation}
which differs from negative Brier loss only by a label-only constant. It is strictly
proper and bounded, avoids probability clipping, and uses the entire predictive
distribution. A floored log score is a sensitivity analysis.

\subsection{Within-cell counterfactual transport}\label{sec:transport}

For a cell $c$, define three population quantities:
\begin{align}
\Delta T_c &= \mathbb E[H_X(Y)\mid C=c], \\
\Delta P_c &= \mathbb E_{H\mid c}\,\mathbb E_{Y'\mid c}[H_X(Y')], \\
\Delta R_c &= \Delta T_c-\Delta P_c .
\end{align}
Here $Y'$ is an independent label drawn from the same cell. Thus $\Delta T_c$ is the
observed score gain, $\Delta P_c$ is the gain retained after breaking the
instance-specific prediction--label assignment while preserving the cell, and
$\Delta R_c$ is the covariance gain lost by that transport.

\begin{theorem}[Transport and covariance identities]\label{thm:population}
For any integrable score contrast $H$,
\begin{equation}
\Delta T_c=\Delta P_c+\Delta R_c,
\qquad
\Delta R_c=\sum_{y=1}^K
\operatorname{Cov}\!\left(H_X(y),\ind \{Y=y\}\mid C=c\right).
\label{eq:population-id}
\end{equation}
For the quadratic score, writing $\Delta q_y(X)=q_1(y\mid X)-q_0(y\mid X)$,
\begin{equation}
\Delta R_c=2\sum_{y=1}^K
\operatorname{Cov}\!\left(\Delta q_y(X),\ind \{Y=y\}\mid C=c\right).
\label{eq:cov-id}
\end{equation}
Consequently, if the gain vector is a function of $c$ alone, then $\Delta R_c=0$.
\end{theorem}

Equation~\eqref{eq:cov-id} is the operational meaning of the construction: the new model
receives covariance credit only to the extent that its probability change aligns more
strongly with the label of the correct case within the same cell. Cell-constant score
improvements, including a better fit to the cell's label histogram, remain in $\Delta P$. A
nonzero covariance does not assert causal or compositional understanding; unrecorded
within-cell shortcuts contribute to it too, so $\phi$ must be declared and stress-tested.

\subsection{What the transported gain measures}\label{sec:transported}

The surviving component is easy to misname. It is tempting to call $\Delta P$ a prior-fit
gain, since transport replaces a case's label by an independent draw from its cell; earlier
drafts of this paper did. Proposition~\ref{prop:transported} in
Appendix~\ref{app:transported} shows the name is wrong and says what the quantity is
instead: under the quadratic score $\Delta P_c$ equals the improvement in how closely the new
model's \emph{mean} within-cell forecast sits to the cell's label frequencies, minus the
increase in the within-cell dispersion of its forecasts. Only the first term is prior fit.
The second rewards the new model for making less dispersed predictions at a fixed mean
forecast, so $\Delta P$ can absorb an entire observed gain with both models' mean forecasts
sitting exactly on the cell frequencies --- a four-point counterexample in that appendix does
just that. This is the reliability-and-sharpness channel of the classical partition, not a
prior-fitting channel, which is why we name it for the operation that produces it; it is also
why confidence matching moves $\widehat{\Delta P}$ so sharply, since a temperature change is
almost purely a change in dispersion.

\subsection{Generalization: the Bregman-score family}\label{sec:bregman}

Nothing in the proof of Theorem~\ref{thm:population} uses properness: the identity and the
covariance sum follow from the definitions alone, for any score. What is specific to the
quadratic score is the further reduction to Eq.~\eqref{eq:cov-id}, in which the part of
$H_x(y)$ not depending on $y$ integrates out because $\sum_y\ind\{Y=y\}=1$ almost surely.
The same cancellation occurs for every strictly proper score generated by a Bregman
divergence, so Eq.~\eqref{eq:cov-id} and the log-score identity used as a robustness check
in Section~\ref{sec:sggaudit} are two instances of one theorem. For a strictly convex,
differentiable $\psi$ on the simplex, the Bregman score is
$S_\psi(q,y)=\psi(q)+\langle\nabla\psi(q),e_y-q\rangle$; taking $\psi(q)=\lVert q\rVert_2^2$
recovers the quadratic score and $\psi(q)=\sum_kq_k\log q_k$ the log score
\citep{gneiting2007proper}.

\begin{theorem}[Bregman-score covariance identity]\label{thm:bregman}
Let $H_x(y)=S_\psi(q_1(\cdot\mid x),y)-S_\psi(q_0(\cdot\mid x),y)$ for a Bregman score generated by a strictly convex, differentiable $\psi$, and write $\Delta g_y(x)=\partial_y\psi(q_1(\cdot\mid x))-\partial_y\psi(q_0(\cdot\mid x))$ for the contrast of the $y$-th gradient coordinate. Then
\begin{equation}
\Delta R_c=\sum_{y=1}^K\operatorname{Cov}\!\bigl(\Delta g_y(X),\,\ind\{Y=y\}\mid C=c\bigr).
\label{eq:bregman-cov-id}
\end{equation}
\end{theorem}

\begin{corollary}[Quadratic and log score]\label{cor:bregman}
For $\psi(q)=\lVert q\rVert_2^2$, $\Delta g_y(x)=2\Delta q_y(x)$ and Eq.~\eqref{eq:bregman-cov-id} is Eq.~\eqref{eq:cov-id}. For $\psi(q)=\sum_kq_k\log q_k$, $\Delta g_y(x)=\log q_1(y\mid x)-\log q_0(y\mid x)$, and
\begin{equation}
\Delta R_c=\sum_{y=1}^K\operatorname{Cov}\!\left(\log\frac{q_1(y\mid X)}{q_0(y\mid X)},\ \ind\{Y=y\}\ \middle|\ C=c\right).
\label{eq:log-cov-id}
\end{equation}
\end{corollary}

Eq.~\eqref{eq:log-cov-id} gives the log-score check of Section~\ref{sec:sggaudit} the same exact covariance meaning that Eq.~\eqref{eq:cov-id} gives the quadratic score: within-group covariance under the log score is the within-cell covariance between the log-likelihood-ratio contrast of the two models and the observed label indicator, not merely "a sensitivity analysis." Any other Bregman score -- the spherical or power-family scores catalogued by \citet{gneiting2007proper}, for instance -- inherits an identity of the same form through its own $\psi$, so Theorem~\ref{thm:bregman} fixes the scope of the estimator's scoring-rule dependence once, rather than leaving each alternative score to be checked separately. The implemented log score floors each probability at $\epsilon=10^{-6}$ to avoid $-\infty$ on an unseen label; Eq.~\eqref{eq:log-cov-id} is exact for the unfloored $\log q_y$, and the floored version used throughout is a sensitivity check on that exact identity rather than a literal instance of it, negligibly different whenever no relevant probability is within $\epsilon$ of zero.

\subsection{Coarsening the cell}\label{sec:coarsening}

Appendix~\ref{sec:ablations} reports that replacing the class-pair cell with a coarser
subject-only cell changes the estimated covariance gain. The following result gives that
sensitivity an exact population-level mechanism rather than treating it as an empirical
curiosity.

\begin{proposition}[Coarsening decomposition]\label{prop:coarsen}
Let $\bar\phi=g(\phi)$ be any coarsening of the grouping variable, so that each cell
$\bar c$ of $\bar\phi$ is a union of cells of $\phi$. Then
\begin{equation}
\Delta R_{\bar c}
=\mathbb E\!\left[\Delta R_C \mid \bar\phi=\bar c\right]
+\sum_{y=1}^K\operatorname{Cov}\!\left(\mathbb E[H_X(y)\mid C],\,
\Pr(Y=y\mid C)\ \middle|\ \bar\phi=\bar c\right).
\label{eq:coarsen}
\end{equation}
\end{proposition}

The first term is the coarse cell's share of the fine-grained covariance gain; the second is
a between-cell covariance, over the finer cells nested inside $\bar c$, between each fine
cell's typical gain at label $y$ and its label-$y$ frequency. It is not sign definite, so
coarsening $\phi$ can inflate or deflate the credited covariance gain, and it vanishes
exactly when either quantity is constant within each coarse cell. This is why the finest
defensible $\phi$ is the conservative default: coarsening adds cross-cell prior structure
into $\Delta R$, while the finer partition's gain is preserved on average through the first
term. The proof is a per-label application of the law of total covariance
(Appendix~\ref{app:proofs}).

\subsection{Sensitivity to an unrecorded confounder}\label{sec:sensitivity}

Read in one direction, Proposition~\ref{prop:coarsen} says what folding a covariate into
$\phi$ does. Read in the other, it quantifies what happens when a relevant covariate is
\emph{never} declared, which is the estimator's main scope limit. Let $Z$ be unrecorded and
let $(\phi,Z)$ be the finer, fully adjusted cell. Applying
Proposition~\ref{prop:coarsen} with $\bar\phi=\phi$ gives, for every declared cell $c$,
\begin{equation}
\Delta R_c(\phi)=\mathbb E\!\left[\Delta R_{(\phi,Z)}\mid\phi=c\right]
+\sum_{y=1}^K\operatorname{Cov}\!\left(\mathbb E[H_X(y)\mid\phi,Z],\,
\Pr(Y=y\mid\phi,Z)\ \middle|\ \phi=c\right),
\label{eq:sensitivity-exact}
\end{equation}
whose first term is what an analyst who could observe $Z$ would report and whose second is
the bias of using $\phi$ alone. The identity is exact but needs $Z$'s joint law with
$(H,Y)$. The following bound replaces that with a judgement about how strongly an unrecorded
variable could plausibly act, in the spirit of \citet{rosenbaum2002observational}.

Proposition~\ref{prop:sensitivity} in Appendix~\ref{app:sensitivity} states the bound: for a
single elicited correlation $\bar\rho$ between an unrecorded variable's effect on the
within-cell gain and its effect on the within-cell label distribution, the bias of using
$\phi$ alone is at most $\bar\rho$ times a product of two dispersion factors, with a
data-free form $\bar\rho M\sqrt{K-1}$ needing only the label cardinality and the score
bound. Because $\bar\rho$ is a correlation rather than a domain-specific odds ratio, it can
be elicited the same way across benchmarks. Appendix~\ref{app:sensitivity} evaluates both
forms on our own estimates, and the sharp one is the less comfortable: an unrecorded
covariate correlated at $0.061$ would account for the entire covariance gain reported in
Section~\ref{sec:results}.

\subsection{Exact finite-sample decomposition}\label{sec:finite}

Consider cell $c$ with $n_c\ge2$ examples. Write $H_i(y)=H_{x_i}(y)$. WAGER computes
\begin{align}
\widehat{\Delta T}_c
 &=\frac1{n_c}\sum_{i\in c} H_i(y_i), \\
\widehat{\Delta P}_c
 &=\frac1{n_c(n_c-1)}\sum_{i\ne j\in c}H_i(y_j),
\label{eq:sample-prior}\\
\widehat{\Delta R}_c
 &=\binom{n_c}{2}^{-1}\sum_{i<j\in c} A_{ij},
\label{eq:sample-reason}
\end{align}
where the antisymmetric kernel is
\begin{equation}
A_{ij}=\frac12\{H_i(y_i)+H_j(y_j)-H_i(y_j)-H_j(y_i)\}.
\label{eq:kernel}
\end{equation}

\begin{theorem}[Finite-sample WAGER identity]\label{thm:sample}
For every cell with $n_c\ge2$, without a distributional assumption,
\begin{equation}
\widehat{\Delta T}_c=\widehat{\Delta P}_c+\widehat{\Delta R}_c.
\end{equation}
Moreover, when the cell's examples are drawn i.i.d.\ from the cell population,
$\widehat{\Delta R}_c$ is an unbiased order-two U-statistic for $\Delta R_c$.
Independence, not exchangeability alone, is what unbiasedness requires: the
independent-coupling definition of $\Delta P_c$ pairs an example with a label drawn
independently from the same cell, so $\mathbb E[H_1(Y_2)]$ matches it only for
independent within-cell draws.
\end{theorem}

Dataset-level quantities weight cells by their identified counts,
$\widehat{\Delta Z}=N_*^{-1}\sum_{c:n_c\ge2}n_c\widehat{\Delta Z}_c$ for $Z\in\{T,P,R\}$.
Singleton cells cannot identify a within-cell counterfactual, so we exclude and report them
rather than inventing a smoothed estimate. If $\widehat{\Delta T}>0$ the descriptive share
$\widehat{\Delta R}/\widehat{\Delta T}$ may exceed one, when covariance improves while the
transported channel deteriorates, and it is undefined for a nonpositive total; we therefore
report $\widehat{\Delta R}$ as the primary result rather than a thresholded ratio.

\subsection{Linear-time computation}\label{sec:efficient}

Naively evaluating Eq.~\eqref{eq:sample-prior} is quadratic in cell size. Let $m_{cy}$
be the count of label $y$ in cell $c$. For example $i\in c$, its transported score is
\begin{equation}
P_i=\frac{\sum_{y=1}^K m_{cy}H_i(y)-H_i(y_i)}{n_c-1}.
\label{eq:linear}
\end{equation}
Averaging $P_i$ yields Eq.~\eqref{eq:sample-prior}; averaging
$H_i(y_i)-P_i$ yields Eq.~\eqref{eq:sample-reason}. Thus the entire estimator costs
$O(NK)$ time and $O(NK)$ storage for cached prediction vectors. No model fitting,
projection fold, hierarchy, or smoothing hyperparameter appears.

\subsection{Why leave-one-out transport, not a fitted-in-sample prior}\label{sec:attenuation}

Equation~\eqref{eq:linear} excludes example $i$'s own label from its transported score.
It is natural to ask what is lost by the simpler alternative of fitting the cell's
empirical label frequency on all $n_c$ examples, including $i$, and using that fitted
frequency as the counterfactual, i.e.\ treating $\hat p_c(y)=m_{cy}/n_c$ as a prior model
and setting $\tilde P_i=\sum_y \hat p_c(y)H_i(y)=n_c^{-1}\sum_y m_{cy}H_i(y)$. This is
the natural plug-in estimator, and it is what a frequency-collapsed projection baseline
computes when no separate evaluation fold is withheld from fitting.

\begin{proposition}[Exact attenuation of the in-sample plug-in]\label{prop:attenuate}
For every cell with $n_c\ge2$ and every $i\in c$, writing $\hat P_i$ for
Eq.~\eqref{eq:linear} and $\hat R_i=H_i(y_i)-\hat P_i$,
\begin{equation}
\tilde P_i=\hat P_i+\frac1{n_c}\hat R_i,
\qquad
\tilde R_i:=H_i(y_i)-\tilde P_i=\frac{n_c-1}{n_c}\,\hat R_i.
\label{eq:attenuation}
\end{equation}
Consequently the in-sample plug-in's dataset-level statistic obeys
$\widetilde{\Delta R}=N_*^{-1}\sum_{c:n_c\ge2}(n_c-1)\widehat{\Delta R}_c$, a
multiplicative shrinkage of every cell's leave-one-out covariance gain by the factor
$(n_c-1)/n_c$.
\end{proposition}

The shrinkage is deterministic, not a variance statement: whenever
$\widehat{\Delta R}_c\neq0$ the in-sample plug-in is biased for it at every finite $n_c$,
worst in the smallest identified cells ($50\%$ at $n_c=2$). The standard remedy for such
self-prediction bias is a projection/evaluation split, which restores unbiasedness by
discarding examples from both stages and introduces a split-fraction hyperparameter and
typically higher variance. Leave-one-out transport removes the bias exactly, at the full
sample, for every cell with $n_c\ge2$, with no fitted nuisance distribution at all.

\subsection{Inference}\label{sec:inference}

For point estimation we use all unordered within-cell pairs. For uncertainty, define
$\bar A_i=(n_c-1)^{-1}\sum_{j\ne i}A_{ij}$ and
$\widehat{\Delta R}_c=n_c^{-1}\sum_i\bar A_i$. The estimated H\'ajek influence
contribution for random cell membership is
\begin{equation}
\widehat\psi_i=2\bar A_i-\widehat{\Delta R}_c-\widehat{\Delta R}.
\label{eq:influence}
\end{equation}
We sum $\widehat\psi_i$ within image and use the resulting one-way sandwich variance,
so multiple relations from the same image are not treated as independent. Total-gain
influence is $H_i(y_i)-\widehat{\Delta T}$; prior-gain and share intervals follow from
the exact identity and delta method.

As a complementary finite-group diagnostic, WAGER applies an independently sampled
cyclic shift to labels inside every cell. Each shift preserves cell size and label
histogram. Recomputing $\widehat{\Delta R}$ under $B$ shifts gives
\begin{equation}
p=\frac{1+\sum_{b=1}^B\ind \{\widehat{\Delta R}^{(b)}\ge
\widehat{\Delta R}\}}{B+1}.
\end{equation}
This is exact under the sharp null that labels are exchangeable relative to the frozen
model outputs within each cell. Because SGG relations share images, our primary
real-data uncertainty statement is the image-cluster-robust interval; the randomization
test is reported as a secondary relation-level diagnostic.

\subsection{Asymptotic normality}\label{sec:asymptotic}

The influence function and sandwich variance behind the reported intervals are derived in
Appendix~\ref{app:influence} through a H\'ajek-projection expansion. The following supplies
the limit theorem, under conditions that hold whenever the score is bounded and no cluster
dominates the sample.

Theorem~\ref{thm:clt} in Appendix~\ref{app:clt} states it: if the score is bounded, examples
fall in mutually independent clusters none of which dominates the sample, the number of
eligible cells is negligible against the identified sample, and $\phi$ has fixed finite
support whose every cell grows, then $\sqrt{N_*}(\widehat{\Delta R}-\theta)$ is
asymptotically normal around the identified-subpopulation estimand and the implemented
sandwich variance is consistent for its limiting variance. The last condition is an
idealisation no heavy-tailed benchmark attains in sample, which is why
Section~\ref{sec:simulation}'s coverage simulation, not this theorem, is what licenses
reading the reported intervals at finite $N$.

\paragraph{Relation to comparative forecast evaluation.}
\begin{corollary}[Relation to Diebold--Mariano/Giacomini--White]\label{cor:dm}
Under the trivial partition $\phi\equiv\text{const}$ (one cell containing every
example), $\widehat{\Delta T}=N^{-1}\sum_iH_i(y_i)$ is the sample mean of the paired loss
differential between the two models, and Theorem~\ref{thm:clt} applied to $\widehat{\Delta T}$
with its image-clustered interval is the grouped, cluster-robust form of the
Diebold--Mariano statistic for that differential \citep{diebold1995comparing}; because
the conditioning information set is empty, it is equivalently a degenerate instance of
the Giacomini--White conditional predictive-ability test \citep{giacomini2006tests}.
\end{corollary}
The corollary is immediate from the definitions: DM and GW test whether a paired loss
differential's mean is zero, and $\widehat{\Delta T}$ at the trivial partition is exactly
that differential's sample mean. Neither test forms $\widehat{\Delta P}$ or
$\widehat{\Delta R}$; both target $\Delta T$ and stop. One qualification on how far the
reading travels: the corollary is stated at the trivial partition, whereas the totals
reported at a nontrivial $\phi$ are $N_*$-weighted means over the \emph{identified}
subsample. Because $\Delta T$ does not depend on $\phi$ except through that exclusion, each
reported $\widehat{\Delta T}$ is a clustered DM statistic for the paired differential on its
identified subsample, which at the coverages here ($99.0\%$ on VG150, $98.9\%$ in
Section~\ref{sec:sggaudit}) differs from the full-sample statistic only in the excluded
singletons.

